% Proved after four complete actual-file audits. The original detecting curve is fixed once; all cover-dependent costs are displayed.
\subsection{Hereditary projective growth from one fixed detecting curve}

\begin{lemma}[Proved: a numerical bound for fixed local projective types on a curve]
\label{D:curve-fixed-local-growth-numerical}
There is a constant \(c_p\), depending only on the positive
characteristic \(p\), with the following property.
Let \(X\) be a smooth proper connected curve of genus \(g\)
over an algebraically closed field of characteristic \(p\),
let \(P\ne\varnothing\) be a finite reduced divisor, and put
\(X^\circ=X-P\).
For \(q=p^f\), fix full local projective representation
classes \(\sigma_x\) in \(\operatorname{PGL}_r(\mathbb F_q)\)
at the points \(x\in P\).
Suppose that the finite Galois splitting extension of
\(\sigma_x\) has normalized different at most \(d_x\ge0\):
its different exponent divided by its ramification index
is at most \(d_x\). Use zero for the trivial extension.
Then the number of global projective conjugacy classes
with these full local types is at most
\[
 q^{\,c_p\left(1+g+|P|+\sum_{x\in P}d_x\right)r^2}.
                                                        \tag{D.GH.1}
\]
The constant is independent of the curve, the local extensions,
their degrees, and the field containing their coefficients.
\end{lemma}

\begin{proof}
We make quantitative the actual atlas and local comparison
in Theorem~\ref{D:curve-fixed-local-linear-growth}.
The affine curve \(X^\circ\) has
\(H^2(X^\circ,\mu_{q-1})=0\), so every projective
representation lifts to \(\operatorname{GL}_r(\mathbb F_q)\)
over the same field, exactly as in
Corollary~\ref{D:curve-fixed-local-projective-growth}.
If the counted fibre is nonempty, one such global lift
supplies one lift of each local type. At each point all
other linear lifts differ by a character of
\(k((t_x))^\times/k((t_x))^{\times(q-1)}\); there are
exactly \(q-1\) such characters. For \(q=2\) there is no
scalar choice.

We need a normalized-different assertion, rather than an
unjustified inference about a highest upper break.
Let \(L/K\) and \(\widetilde L/K\) be the splitting
extensions of a projective local type and any linear lift.
Then \(L\subset\widetilde L\), and the group of
\(\widetilde L/L\) is a subgroup of
\(\mathbb F_q^\times\). It is therefore tame.
All these complete fields have the same algebraically
closed residue field. If \(e=e(L/K)\) and
\(\widetilde e=e(\widetilde L/L)\), transitivity of the
different gives
\[
 \frac{\delta(\widetilde L/K)}{e\widetilde e}
  =\frac{\delta(L/K)}e+
       \frac{\widetilde e-1}{e\widetilde e}
  \le d_x+1.                                           \tag{D.GH.2}
\]
In this equal-characteristic scope the transitivity formula
also follows immediately from the chain rule: for
\(t=F(s)\), \(s=G(w)\), the valuation of
\((F\circ G)'(w)=F'(G(w))G'(w)\) is the sum of the
two different valuations with the indicated ramification
factor. The derivative formula for the different is
\citet[Lemmas 49.12.1--3 and Example 49.12.5,
Tag 0BWB]{StacksDifferentLatticeB}.

Fix one tuple of linear local lifts. Put
\[
 c_x=\left\lceil p(d_x+1)\right\rceil,\qquad
 B=\sum_{x\in P}c_xx,\qquad
 \Delta=\deg B,\qquad
 A=1+g+\Delta+|P| .
                                                        \tag{D.GH.3}
\]
Lemma~\ref{D:different-Frobenius-invariant-lattice}
gives the canonical invariant lattices for all coefficient
phases, with forward Frobenius maps holomorphic and their
inverse maps of pole at most \(B\).
Thus the canonical bundles \(E_j\), \(j\in\mathbb Z/f\),
on \(X\) have actual arrows
\[
 F_X^*E_j\longrightarrow E_{j+1},
 \qquad E_{j+1}(-B)\longrightarrow F_X^*E_j,
                                                        \tag{D.GH.4}
\]
with both inverse products the Cartier inclusion.
The slope argument of (D.CL.7), with \(B\) in place of
\(cD\), gives
\[
       -\frac{\Delta}{p-1}\le\mu_{\min}(E_j)
                  \le\mu_{\max}(E_j)\le0.               \tag{D.GH.5}
\]
This argument uses the two actual injections and their
duals; it assumes no tensor-product semistability.

Here are choices of all presentation degrees.
Choose a line bundle \(H\) of degree
\[
 h=2g+4+\left\lceil\frac{\Delta}{p-1}\right\rceil
       =O_p(A).
                                                        \tag{D.GH.6}
\]
Such a line exists, for example as a divisor supported
at one rational point. The bundle \(E_j^+=E_j\otimes H\)
has minimum slope at least \(2g+4\), and hence has
vanishing \(H^1\) and is generated. A general
\((r+1)\)-tuple of sections generates it, with line
kernel \(L_j^{-1}\), where \(L_j=\det E_j^+\).
The degree bounds are
\[
       r(2g+4)\le\deg L_j\le rh.
                                                        \tag{D.GH.7}
\]
Consequently both \(L_j\) and \(L_j^{-1}H^{3r}\)
are generated and have vanishing \(H^1\).
The fixed partial-resolution construction used in
(D.CL.8) therefore applies with
\[
 P_0=\mathcal O_X^{r+1},\qquad
 P_i=(H^{-3ir})^{\oplus2}\quad(1\le i\le7).             \tag{D.GH.8}
\]
The alternating line syzygies are exactly those in that
construction. All differential coefficients are
parameters in spaces of global sections of these fixed
terms; no moving determinant or variable cohomology
basis is introduced.

Impose \(d^2=0\) and the actual open rank conditions.
Local unit-minor cancellation makes the degree-zero
cohomology a bundle over every coefficient algebra.
Use length seven for source complexes, length three for
targets, and the same top-generator identity test as in
Theorem~\ref{D:curve-fixed-local-linear-growth}.
The two end-error intervals are disjoint from the Hom
degrees used for maps and identities. Thus the finite
Hom cycles and degree-minus-one boundary witnesses
express precisely (D.GH.4) and their inverse products.

We verify the dependence on the changing curve.
For any line of degree \(d\) on \(X\),
\(h^0+h^1\le2g+2+|d|\); this follows from a regular
section when one exists and Riemann--Roch, and is
immediate when there is no section.
Every top--top Hom block has multiplicity \(O(r^2)\)
and line degree \(O_p(h+\Delta)\). A block with one
positive-index term has multiplicity \(O(r)\) and
line degree \(O_p(rh+\Delta)\); a block with two
positive-index terms has bounded multiplicity and
the same degree bound. The twists \(H\), \(B\), and
one-step Frobenius in (D.GH.4) only change the
absolute constants depending on \(p\). Hence each
phase has
\[
                   O_p(A r^2)                         \tag{D.GH.9}
\]
object and Hom coordinates.

As in the proved fixed-type atlas, split the two-affine
Cech complexes of the fixed line terms over the ground
field. The internal widths are bounded by seven and by
the fixed map-composition diagram. The finite homological
perturbation identities therefore use a bounded number
of factors, independent of \(g,\Delta,r,f\).
They commute with every coefficient-base change and
give polynomial equations of degree bounded by a constant
depending only on \(p\). The only coefficient power in
one phase is the \(p\)-th power. No expanded \(q\)-Frobenius
identity occurs.

For each \(x\), choose
\[
 N_x=\left\lfloor\frac{c_x}{p-1}\right\rfloor+1,
 \qquad L_x=N_x+c_x .
                                                        \tag{D.GH.10}
\]
Then \((p-1)N_x>c_x\) and
\(\sum_xL_x=O_p(\Delta+|P|)\).
Frame the quotients of (D.GH.8) on
\(\coprod_x\operatorname{Spec}k[t_x]/(t_x^{L_x})\)
using the explicit matrices \(u,v,w\) of (D.CL.9),
and impose the finite cyclic unit equations with the
fixed canonical local models. These add
\[
                  O_p\!\left(fr^2\sum_xL_x\right)
                         =O_p(f A r^2)                \tag{D.GH.11}
\]
variables and equations of bounded degree.
Increasing the truncation length increases the number
of coefficient variables, not the polynomial degree.
Lemma~\ref{D:curve-integral-unit-jet-comparison}
gives the exact integral local identifications over
all coefficient algebras.

Every geometric point of the resulting polynomial
zero scheme and open rank locus is an actual global
unit bundle with these fixed canonical local types,
and every representation counted is covered.
The entire specified-Artinian Isom proof in
Theorem~\ref{D:curve-fixed-local-linear-growth}
applies unchanged: a unit isomorphism lifts uniquely
on \(X^\circ\), its local lift preserves the fixed
canonical invariant lattices, and faithful flatness
of completion extends it over the proper curve.
The finite-presentation Isom is etale and its distinct
class images are disjoint open subsets. This uses
the fixed full local types, not just the different
bounds.

There are at most \(C_p f A r^2\) affine variables
and the equation degrees are at most \(D_p\), with
constants depending only on \(p\).
The component bound, including its prefactor, is
\[
 (C_p f A r^2+1)\max\{2,D_p\}^{C_p f A r^2}
                         \le p^{C'_p f A r^2}.         \tag{D.GH.12}
\]
Passing to an open locus does not add irreducible
components. Finally multiply by the at most
\((q-1)^{|P|}\) tuples of linear local lifts.
Since \(\Delta\le p\sum_xd_x+(p+1)|P|\),
this proves (D.GH.1). The case of an empty counted
fibre, and the trivial projective rank-one case,
are included without any additional choices.
\end{proof}

\begin{theorem}[Proved: global hereditary growth with bounded local periods]
\label{D:global-finite-puncture-hereditary-bounded-local-period}
Let \(S\) be a fixed normal proper connected surface
over an algebraically closed field \(k\) of characteristic
\(p>0\), let \(T\subset S\) be finite, and put
\(\Gamma=\pi_1(U)\), \(U=S-T\).
Choose the fixed finite set \(F\supset T\), complete
local sources \(R_s\), and detecting curve of
Lemma~\ref{D:finite-puncture-detecting-curve-local-types}.
Fix \(B_*\ge1\).

For an open \(H\subset\Gamma\) of index \(t\), let
\(U_H\to U\) be its connected finite etale cover.
For each \(s\in F\), write the pullback to the fixed
puncture \(U_{R_s}\) as its connected components:
\[
 U_H\times_U U_{R_s}=\coprod_i U_{s,i},
 \qquad n_{s,i}=\deg(U_{s,i}/U_{R_s}),\qquad
                       \sum_i n_{s,i}=t.               \tag{D.GH.13}
\]
Each \(U_{s,i}\) is the puncture of its complete
normal local normalization of \(R_s\).
There is a constant \(A_{S,T,B_*}\) such that the
homomorphisms
\(\rho:H\to\operatorname{PGL}_r(\mathbb F_q)\),
\(q=p^f\), for which every restriction to
\(\pi_1(U_{s,i})\) has geometric Brauer period at most
\(B_*\), have at most
\[
                       q^{A_{S,T,B_*}t r^2}             \tag{D.GH.14}
\]
projective conjugacy classes.

It suffices to require the period condition at
\(s\in T\); the other restrictions are trivial.
In particular the conclusion covers every such
representation if \(\gcd(r,q-1)\le B_*\).
A bound \(B_*\) on the global geometric Brauer
period on \(U_H\) is another sufficient condition.
The detecting curve is chosen once for \((S,T)\).
No constant belonging to a new unspecified curve
or to a new surface resolution is used.
\end{theorem}

\begin{proof}
Write \(C^\circ=C-P\) for that one fixed detecting
curve over its fixed algebraically closed extension
\(\Omega/k\), and let \(b_x\) be its fixed upper-break
bounds at \(x\in P\).
The pullback \(C_H^\circ=C^\circ\times_U U_H\)
is connected, since
\(\pi_1(C^\circ)\twoheadrightarrow\Gamma\).
Let \(C_H\) be its smooth proper completion and
\(P_H=C_H-C_H^\circ\). Then \(C_H\to C\) is a
finite separable morphism of degree \(t\), and
\(\pi_1(C_H^\circ)\twoheadrightarrow H\).
All these are pullbacks of the original detector;
no further Bertini choice is made.

For \(y\in P_H\) lying over \(x\in P\), put
\[
 K_x=\Omega((t_x)),\quad K_y=\Omega((s_y)),\quad
 e_y=[K_y:K_x],\quad \delta_y=\delta(K_y/K_x).
                                                        \tag{D.GH.15}
\]
The residue fields are \(\Omega\), so the displayed
degree is the ramification index and
\(\sum_{y\mid x}e_y=t\).
Choose a finite Galois quotient of \(\Gamma\) whose
kernel is contained in \(H\). Its local splitting
field at \(x\) contains every \(K_y\) after compatible
embeddings. The original detecting-curve lemma bounds
its normalized different by \(b_x+1\).
The tower formula (D.GH.2), without its tame
specialization, implies
\[
             0\le\delta_y\le e_y(b_x+1).
                                                        \tag{D.GH.16}
\]
The Riemann--Hurwitz formula with the actual different
therefore gives
\[
 \begin{aligned}
 2g(C_H)-2
   &=t(2g(C)-2)+\sum_{y\in P_H}\delta_y,\\
 |P_H|&\le t|P|,\qquad g(C_H)+|P_H|+1=O_{S,T}(t).
 \end{aligned}                                         \tag{D.GH.17}
\]
The exact wild different formula and its derivation
are \citet[Section 53.12 and Lemma 53.12.2,
Tag 0C1B]{StacksResidualGenusB}. We have not used the
tame replacement \(\delta_y=e_y-1\).

The same argument also controls all representations
of \(H\), rather than only the degree-\(t\) cover.
Given any finite representation of \(H\), its kernel
has an open normal core in \(\Gamma\).
The corresponding finite Galois quotient of
\(\Gamma\) simultaneously trivializes that
representation and the coset cover.
Let \(L/K_x\) be its local splitting field.
It is Galois, contains \(K_y\), and satisfies
\(\delta(L/K_x)/e(L/K_x)\le b_x+1\).
Transitivity gives the exact equality and bound
\[
 \frac{\delta(L/K_y)}{e(L/K_y)}
 =e_y\,\frac{\delta(L/K_x)}{e(L/K_x)}-\delta_y
 \le d_y,\qquad
                  d_y=e_y(b_x+1)-\delta_y\ge0.          \tag{D.GH.18}
\]
Passing to a Galois quotient over \(K_y\) can only
decrease its normalized different, again by
transitivity. Thus \(d_y\) bounds the projective
local splitting type of every representation
under consideration.
In particular
\[
                 \sum_{y\in P_H}d_y
                      \le t\sum_{x\in P}(b_x+1)
                      =O_{S,T}(t).                    \tag{D.GH.19}
\]
For any same-field linear lift on \(C_H^\circ\),
the scalar tame extension adds at most one to
each normalized different by (D.GH.2).
Consequently the local inverse-Frobenius depths
\(\lceil p(d_y+1)\rceil\), and the retained and
testing jet lengths (D.GH.10), have total \(O_{S,T,p}(t)\).
Using the maximum local depth at all marked points
would lose this estimate and is unnecessary.

It remains to count the possible tuples of full local
projective types. Each source \(R_s\) is fixed before
\(H\) is chosen. Its completed finite cover factors
in (D.GH.13) have total degree \(t\), even when
the maps of local fundamental groups to \(\Gamma\)
are not injective.
Theorem~\ref{D:bounded-period-hereditary-projective-growth}
therefore bounds all representation classes on
these factors by
\[
 \prod_{s\in F}\prod_i q^{A_{R_s,B_*}n_{s,i}r^2}
                  =q^{t r^2\sum_{s\in F}A_{R_s,B_*}}.
                                                        \tag{D.GH.20}
\]
The factorization of the original formal branches
through \(U_{R_s}\) pulls back to factorizations
through the appropriate \(U_{s,i}\).
The finite etale algebra over \(K_x\) is the product
of the fields \(K_y\), so these are precisely all
the full local types on \(C_H\).
Choose compatible paths once for this fixed cover.
One class on a factor \(U_{s,i}\) determines, by
restriction, the classes at every curve branch
using that factor. Hence (D.GH.20) counts each
surface factor once. There is no extra factor
equal to the number of branches or the local
cover degree. Independent choices of different
surface factors can overcount, which is harmless.

For each retained tuple apply
Lemma~\ref{D:curve-fixed-local-growth-numerical}
to \(C_H,P_H\), with the \(d_y\) of (D.GH.18).
Equations (D.GH.17)--(D.GH.19) bound its exponent
by a constant times \(t r^2\).
Multiplying by (D.GH.20) proves the same bound
for the corresponding projective classes on
\(C_H^\circ\). Precomposition with the surjection
\(\pi_1(C_H^\circ)\twoheadrightarrow H\) is injective
on conjugacy classes, proving (D.GH.14).

If \(s\notin T\), every original cover of \(U\)
is split over \(\operatorname{Spec}R_s\), as in
the detecting-curve lemma; the induced local
representations are trivial. A global geometric
period bound restricts to every factor in
(D.GH.13). Finally the ordinary projective
multiplier is killed by both \(r\), by taking
determinants of linear lifts, and \(q-1\), by
its coefficient exponent. Its geometric image
therefore has period dividing \(\gcd(r,q-1)\).
These prove the stated sufficient conditions.
\end{proof}

\begin{remark}[Scope of the hereditary transfer]
The new uniformity concerns arbitrary finite covers of
one fixed finite-puncture proper surface and fixed
bounded local geometric periods. It is stronger than
reapplying a fixed-surface theorem separately upstairs:
the genus, total different, local jet length, and
surface-source indices have all been bounded linearly
in the original index. It does not claim unrestricted
local period, a semilinear projective transfer, or
finite generation of the whole global group.
This is a transfer of the same bounded-period
hereditary projective enlargement, not a separate
independent reduction.
\end{remark}
